\documentclass[a4paper, 12pt, oneside, headspline]{scrartcl}
% !TeX root = ../expose.tex
% Variablen welche innerhalb der gesamten Arbeit zur Verfügung stehen sollen
\newcommand{\titleDocument}{Exposé}
\newcommand{\subjectDocument}{Vergleich von Python Test-Tools für Test-driven development}
\newcommand{\institution}{Hochschule für angewandte Wissenschaften Augsburg}
\newcommand{\insSection}{Fakultät für Informatik}
\newcommand{\name}{Maximilian Konter}
\newcommand{\email}{maximilian.konter@hs-augsburg.de}


% Farbige links
\usepackage{color}

% Umlaute unter UTF8 nutzen
\usepackage[utf8]{inputenc}

% deutsche Silbentrennung
\usepackage[ngerman]{babel}

% Grafiken aus PNG Dateien einbinden
\usepackage{graphicx}

% bricht lange URLs "schoend" um
\usepackage[hyphens,obeyspaces,spaces]{url}

% mehrseitige Tabellen ermöglichen
\usepackage{longtable}

% für Tabellen
\usepackage{array}

% Schaltet den zusätzlichen Zwischenraum ab, den LaTeX normalerweise nach einem Satzzeichen einfügt.
\frenchspacing

% Paket für Zeilenabstand
\usepackage{setspace}

% für Bildbezeichner
\usepackage{capt-of}

% für Stichwortverzeichnis
\usepackage{makeidx}

% Für Glossar
\usepackage{glossaries}

% Bessere Listen
\usepackage{enumitem}

% Festlegung Art der Zitierung - Havardmethode: Abkuerzung Autor + Jahr
\bibliographystyle{alphadin}

% für Listings
\usepackage{listings}
\lstset{numbers=left, numberstyle=\tiny, numbersep=5pt, keywordstyle=\color{black}\bfseries, stringstyle=\ttfamily,showstringspaces=false,basicstyle=\footnotesize,captionpos=b}
\lstset{language=python}

% Indexerstellung
\makeindex

% Abkürzungsverzeichnis
\usepackage[german]{nomencl}
\let\abbrev\nomenclature

% Abkürzungsverzeichnis LiveTex Version
\renewcommand{\nomname}{Abkürzungsverzeichnis}
\setlength{\nomlabelwidth}{.25\hsize}
\renewcommand{\nomlabel}[1]{#1 \dotfill}
\setlength{\nomitemsep}{-\parsep}
\makenomenclature
\makeglossaries

% Disable single lines at the start of a paragraph (Schusterjungen)
\clubpenalty = 10000
% Disable single lines at the end of a paragraph (Hurenkinder)
\widowpenalty = 10000
\displaywidowpenalty = 10000

% Hyperlinks im Text
\usepackage[bookmarksnumbered,pdftitle={\titleDocument},hyperfootnotes=false]{hyperref}
\hypersetup{colorlinks, citecolor=red, linkcolor=blue, urlcolor=black}

% neue Kopfzeilen mit fancypaket
\usepackage{fancyhdr} %Paket laden
\pagestyle{fancy} %eigener Seitenstil
\fancyhf{} %alle Kopf- und Fußzeilenfelder bereinigen
\fancyhead[L]{\nouppercase{\leftmark}} %Kopfzeile links
\fancyhead[C]{} %zentrierte Kopfzeile
\fancyhead[R]{\thepage} %Kopfzeile rechts
\renewcommand{\headrulewidth}{0.4pt} %obere Trennlinie
\fancyfoot[C]{\thepage} %Seitennummer
\renewcommand{\footrulewidth}{0.4pt} %untere Trennlinie

\begin{document}	
	% Titelseite
	% !TeX root = ../bachelor.tex
\thispagestyle{empty}

\begin{figure}[t]
 \centering
 \includegraphics[width=0.3\textwidth]{abb/logoHSA}
\end{figure}


\begin{verbatim}


\end{verbatim}

\begin{center}
\Large{\institution}
\end{center}


\begin{center}
\Large{\insSection}
\end{center}
\begin{verbatim}




\end{verbatim}
\begin{center}
\doublespacing
\textbf{\LARGE{\titleDocument}}
\singlespacing
\begin{verbatim}

\end{verbatim}
\textbf{{~\subjectDocument~}}
\end{center}
\begin{verbatim}

\end{verbatim}
\begin{center}

\end{center}
\begin{verbatim}

\end{verbatim}
\begin{center}
\textbf{zur Erlangung des akademischen Grades \\ Bachelor of Science}
\end{center}
\begin{verbatim}





\end{verbatim}

\begin{footnotesize}
\begin{flushleft}
\begin{tabular}{llll}
\textbf{Thema:} & & \subjectDocument & \\
& & \\
\textbf{Autor:} & & \name & \\
& & \email & \\
& & MatNr. 951004 & \\
& & \\
\textbf{Version vom:} & & \today &\\
& & \\
\textbf{1. Betreuer:} & & Dipl.-Inf. (FH), Dipl.-De Erich Seifert, MA &\\
\textbf{2. BetreuerIn:} & & Prof. Dr.-Ing. Alexandra Teynor &\\
\end{tabular}
\end{flushleft}
\end{footnotesize}
	
	% Inhaltsverzeichnis anzeigen
	\newpage
	\tableofcontents
	
	% das Abkürzungsverzeichnis
	\newpage
	% Abkürzungsverzeichnis soll im Inhaltsverzeichnis auftauchen
	\addcontentsline{toc}{section}{Abkürzungsverzeichnis}
	% das Abkürzungsverzeichnis entgültige Ausgeben
	\fancyhead[L]{Abkürzungsverzeichnis} %Kopfzeile links
	% !TeX root = ../../main.tex
\nomenclature{GPL}{GNU General Public License}
\nomenclature{GNU}{GNU is not Unix}
\nomenclature{LGPL}{GNU Lesser General Public License}
\nomenclature{TDD}{Test-driven development}
\nomenclature{GUI}{Graphical User Interface}
\nomenclature{STDLIB}{Python Standard Bibliothek}
\nomenclature{VCS}{\gls{vcs}}
	\printnomenclature
    
    % das Glossar
    \newpage
    % Glossar soll im Inhaltsverzeichnis auftauchen
    \addcontentsline{toc}{section}{Glossar}
    % das Abkürzungsverzeichnis entgültige Ausgeben
    \fancyhead[L]{Glossar} %Kopfzeile links
    \newglossaryentry{Pythonista}{
    name={Pythonista},
    description={Jemand, der Python als seine favorisierte Sprache verwendet}
}
\newglossaryentry{Continous testing}{
    name={Continous testing},
    description={Bezeichnet den Prozess des automatischen Testens mit Hilfe
    eines Automatisierungstools. Hierbei werden die Tests durch ein Ereignis
    automatisch ausgeführt, wodurch der Entwickler sofort ein Feedback bekommt}
}
\newglossaryentry{mock}{
    name={mock},
    description={Etwas mocken bedeutet, ein Objekt durch ein falsches Objekt, den
    Mock zu ersetzen, das genau so aussieht, wie das erwartete Objekt, aber nicht das
    Gleiche ist. So ist es möglich eine Funktion zu mocken, wodurch nicht die Funktion,
    sondern der Mock aufgerufen wird und das eingestellte ergebnis liefert}
}
\newglossaryentry{commit}{
    name={commit},
    description={Ein commit ist die Sammlung von änderungen an Dateien, welche 
    mithilfe eines VCS verwaltet werden},
    plural=commits
}
\newglossaryentry{vcs}{
    name={Versions Control System},
    description={Ein VCS ist ein Dienst, der es seinen Nutzern ermöglicht änderungen
    an Dateien in einer Chronik zu speichern um später darauf zu zu greifen. Dies
    dient auch der verteilung von Daten über mehrere Systeme sowie sicherung der
    Daten vor verlust}
}
\newglossaryentry{fuzz}{
    name={fuzz-testing},
    description={Beschreibt eine Art von Test,
    bei dem das zu testende Modul oder Programm mit zufälligen Werten aufgerufen wird. Dabei
    soll jede mögliche anwendung des Moduls oder Programms dargestellt werden}
}
\newglossaryentry{bug}{
    name={bug},
    description={Ein Bug(zu deutsch Käfer) ist ein Fehler in einem Programm oder in einer Software}
}
\newglossaryentry{fixture}{
    name={fixture},
    description={Beschreibt die Präperation von Tests mit festen (fix) Werten, so wird zum Beispiel
    vor jedem Test eine Variable gesetzt. Dies dient der übersichtlichkeit, da dies dann nicht in
    jedem Test vorgenommen werden muss und um varianz zwischen Tests und Test läufen zu vermeiden}
}
\newglossaryentry{cli}{
    name={Command-line interface},
    description={Eine Benutzerschnitstelle für das Terminal}
}
\newglossaryentry{stacktrace}{
    name={StackTrace},
    description={Der StackTrace ist ein Protokoll der aufgerufenen Funktionen und Methoden die 
    ineinander verschachtelt sind, bis hin zu der tiefsten Funkion/Methode die einen Fehler
    geworfen hat}
}
    \printglossaries
	
	%%%%% Beginn Dokument %%%%%
	\newpage
	\fancyhead[L]{\nouppercase{\leftmark}} %Kopfzeile links
	
	% 1,5 facher Zeilenabstand
	\onehalfspacing
	
	% einzelne Kapitel
	% !TeX root = ../expose.tex
\section{Ausgangslage}
TDD wird in der heutigen Softwareentwicklung immer verbreiteter und beliebter.
Die Ansprüche an Software sind in den letzten Jahren immer weiter gestiegen. Dies
liegt vor allem an der Reichweite, die Software heute hat. So besitzen im Jahr
2018 bereits circa 66\% aller Menschen ein Smartphone \cite{FraukeSchobelt:Smartphone},
im Arbeitsleben ist ein PC meist gar nicht mehr weg zu denken. Doch mit den steigenden
Nutzerzahlen steigen auch die Anforderungen, welche die Nutzer an die Software stellen.
Somit wird die Anzahl der gefundenen Bugs dementsprechend größer.

\section{Problembeschreibung}
Im weltweiten Markt gibt es viele große Unternehmen, die gegenseitig um die
Nutzer kämpfen. Selbstverständlich präferieren die Nutzer denjenigen Anbieter, welcher die bessere
Software bietet. Dies kann sich heute jedoch stetig ändern. Mit der steigenden Anzahl
an Bugs, die gefunden werden, steigt auch die Anzahl der Nutzer, die von diesen Bugs betroffen
sind. Diese Bugs sollen natürlich schnellstmöglich gefixt werden um so zu verhindern,
dass die Nutzer die Software wechseln.

So schwer es ist seine Nutzer zu halten, umso schwerer ist es, zum Start einer Software
Nutzer zu akquirieren. Es gibt bereits Software, die ähnliche Services anbieten.
So ist es noch schwerer dem Markt beizutreten. Die Anforderungen werden durch die
bereits am Mark tätigen Firmen gesetzt, damit sollten Fehler, die bereits gelößt
wurden nicht mehr auftauchen.

Für Unternehmen sind diese Anforderungen meist schwer zu meistern, weshalb Software
meist mit Fehlern released wird, um diese dann von den Nutzern aufdecken zu lassen und
zu fixen.

\section{Fragestellung}
Welche Möglichkeiten stehen einem Entwickler/Arbeitgeber zur Verfügung Python
als Sprache für TDD einzusetzen? Daraus resultierend stellt sich die Frage,
welches der Test-Tools das Beste dafür ist, oder welche sich kombinieren lassen
um das Beste aus ihnen heraus zu holen.

Welchen Vorteil zieht ein Entwickle/Arbeitgeber daraus TDD zu verwenden? Ist
der betriebene Aufwand und die daraus resultierende Test-Abdeckung es wert TDD
anderen Techniken vorzuziehen?
	
	% !TeX root = ../expose.tex
\section{Zielsetzung}
Das Ziel dieser Arbeit ist es, Test-Tools der Programmiersprache Python in den
Aspekten Anwendbarkeit, Effizienz, Komplexität und Erweiterbarkeit zu
vergleichen. Daraus resultierend soll eine Empfehlung entstehen, wie man
möglichst einfach TDD mit Python betreiben kann, um so die Motivation zu steigern
Tests zu schreiben.

Um die Effizienz von TDD zu steigern, wird in dieser Arbeit auf die Automatisierung
der Tests eingegangen. Dies wird mit Hilfe des Tools
Travis CI\footnote{\href{https://travis-ci.org/}{Travis CI Website}}
dargestellt.

Die Arbeit soll Anwendern einen Leitfaden yur verfügung stellen welche Tools
zur Zeit zum betreiben von TDD die besten sind.
Nach dem Vergleich soll der Leser ein Bild darüber haben, welche Tools ihm zur
Verfügung stehen um Tests zu schreiben oder TDD zu betreiben. Zusätzlich soll
er ein wenig Verständnis darüber haben, wie er mit Hilfe von
\gls{Continous testing} effektiver TDD betreiben kann.
\newline
Am Ende der Arbeit werden die Vor- und Nachteile von TDD abgewogen, für die
Nachteile soll, soweit dies möglich ist, eine Lösung angeboten werden um TDD
attraktiver zu gestalten.

	% !TeX root = ../expose.tex
\section{Theoretische Grundlage}
Hier aufgeführt findet sich einiges an Literatur, welche relevant für diese Arbeit
sein könnte:

\begin{itemize}
    \item 5-jährige Studie von IBM aus dem Jahr 2007
    \url{http://citeseerx.ist.psu.edu/viewdoc/download?doi=10.1.1.104.6319&rep=rep1&type=pdf}
    \item Diskussion zu einem Buch über Empirische Softwareentwicklung (Paywall)
    \url{https://www.infoq.com/news/2009/03/TDD-Improves-Quality}
    \item Studie über den Gebrauch von TDD und dem daraus resultierenden Design
    und Testverbesserungen
    \url{https://arxiv.org/pdf/1711.05082.pdf}
    \item TDD vs Test-Last
    \url{https://www.researchgate.net/publication/315743099_An_Experimental_Evaluation_of_Test_Driven_Development_vs_Test-Last_Development_with_Industry_Professionals}
    \item TDD vs nicht TDD auf Github
    \url{https://peerj.com/preprints/1920/}
    \item Beispiel für TDD mit Python
    \url{https://medium.freecodecamp.org/learning-to-test-with-python-997ace2d8abe}
    \item Getting started with testing in Python
    \url{https://realpython.com/python-testing/}
    \item Minimum viable testing
    \url{https://realpython.com/the-minimum-viable-test-suite/}
    \item Continuous Integration With Python: An Introduction
    \url{https://realpython.com/python-continuous-integration/}
    \item Test-driven development with Python (Buch)
    \url{https://www.obeythetestinggoat.com/}
\end{itemize}

	% !TeX root = ../expose.tex
\section{Konzept}
Die Test-Tools aus der Standard Bibliothek von Python, sowie die beliebtesten
Test-Tools, werden anhand folgender Aspekte verglichen:

\begin{itemize}
    \item Anwendbarkeit
        \subitem Bietet das Tool alles, um TDD betreiben zu können? (Bei TDD
        kann man Tests \gls{mocken}, wodurch sie nicht fehlschlagen.)
        Mit wie vielen Packeten muss das Tool betrieben werden?
    \item Effizienz
        \subitem Wie viel lässt sich mit diesem Tool möglichst einfach und
        schnell erreichen? Ist besonders viel Vorarbeit notwendig um die Tests
        auf zu setzen oder kann sofort mit dem Schreiben der Tests begonnen
        werden?
        \newline
        \\
        Genauso stellt sich die Frage, wie Effizient der Entwickler die Tests
        auswerten kann.
    \item Komplexität
        \subitem Wie komplex ist das Tool? Das heißt, wie viel Funktionalität
        bietet das Tool dem Entwickler von Haus aus, aber auch wie schwer
        ist es einen Code zu schreiben oder wie schnell wird ein Code unübersichtlich, da
        das Tool viel Code abseits der Tests benötigt.
    \item Erweiterbarkeit
        \subitem Wie leicht lässt sich das Tool mit anderen Tools erweitern?
        Gibt es vielleicht Erweiterungen der Community für dieses Tool, die sehr
        hilfreich sind?
\end{itemize}
Eventuell werden beim Schreiben der Arbeit weitere Aspekte entstehen, die hier
aufgeführt werden können.
\newline
\\
Der Vergleich der Tools entsteht durch das Analysieren dieser auf der
jeweiligen Website, sowie durch ein komplexes Code-Beispiel, das mit dem
jeweiligen Tool geschrieben wird. Passend zum Code-Beispiel wird auch analysiert,
wie einfach die Automatisierung mit Travis CI verläuft, indem dies auch mit einem
echtem Code getestet wird.

Sind die Grundlagen für den Vergleich geschaffen, werden die Tools
untereinander verglichen und ihre Schwächen und Stärken noch zusammen
gefasst.
\newline
\\
Sofern dies möglich ist wird versucht die Tools zu kombinieren um etwaige
Schwächen eines Tools zu auszugleichen.

Nachdem das/die beste(n) Tool(s) gefunden wurde(n), wird auf die Schwächen von TDD
im Allgemeinen eingegangen. Dies behandelt den erforderlichen Aufwand und die
daraus resultierende Zeit, die "verloren" geht (Zeit, die statt für Tests
schreiben auch für Code schreiben genutzt werden könnte), sowie die Disziplin,
die den Entwicklern abverlangt wird, zuerst Tests zu schreiben, bevor die
tatsächliche Funktionalität geschrieben wird. Es wird auch auf den
Wirtschaftlichen Einsatz von TDD eingegangen.

Am Ende der Arbeit nehme ich selbst Stellung zu der Effizienz, Funktionalität und positiven, sowie negativen Aspekten zu TDD. Außerdem gebe ich eine Empfehlung ab, für wen TDD am Besten von Nutzen ist.

	% !TeX root = ../expose.tex
\section{Vorläufige Gliederung}
\begin{enumerate}[label*=\arabic*.]
    \item Abstract
    \item Einleitung
    \begin{enumerate}[label*=\arabic*.]
        \item Was ist Python?
        \item Was ist Test-driven development?
    \end{enumerate}
    \item Python Test-Tools
    \begin{enumerate}[label*=\arabic*.]
        \item Welche Tools bietet die Standard Bibliothek von Python?
        \item Welche Tools gibt es abseits der Standard Bibliothek?
        \begin{enumerate}[label*=\arabic*.]
            \item Python Test-Tools
            \item Python-Framework Test-Tools (Nicht im Vergleich)
            \item Weitere Tools
        \end{enumerate}
        \item Tool XY
        \begin{enumerate}[label*=\arabic*.]
            \item Code-Beispiel
            \item Anwendbarkeit
            \item Effizienz
            \item Komplexität
            \item Erweiterbarkeit
        \end{enumerate}
    \end{enumerate}
    \item Zusammenfassung
    \begin{enumerate}[label*=\arabic*.]
        \item Tool XY
    \end{enumerate}
    \item Vergleich der Tools
    \item Kombinierung von Tools
    \item Diskussion: Lohnt sich TDD
    \begin{enumerate}[label*=\arabic*.]
        \item Welche Stärken hat TDD?
        \item Welche Schwächen hat TDD?
        \begin{enumerate}[label*=\arabic*.]
            \item Welche Möglichkeiten gibt es diese Schwächen zu umgehen?
            \item Welche Schwächen müssen akzeptiert werden?
        \end{enumerate}
        \item Welche wirtschaftlichen Aspekte müssen bei TDD beachtet werden?
        \item Zusammenfassung
    \end{enumerate}
    \item Fazit
    \item Nachwort
\end{enumerate}

	% !TeX root = ../expose.tex
\section{Motivation}
Da ich selbst ein \gls{Pythonista} bin, war es für mich naheliegend meine Bachelorarbeit
dieser Sprache zu widmen. Durch mein Praxissemester und dem daraus resultierenden
Werkstudentenjob, in dem ich seit eineinhalb Jahren an der Test- und Buildautomatisierung
arbeite, wurde mein Interesse in die automatisierte Ausführung von Tests geweckt. Auch mein
Bedürfnis, Tests für meine Programme zu schreiben, ist enorm gewachsen, da ich festgestellt
habe, wie gut das Gefühl ist, wenn man etwas im Code ändert und alle Tests danach noch
zu 100\% durchlaufen.

Seitdem ist Testen ein Teil meiner Entwicklung und auch meines Denkens. TDD selbst habe
ich jedoch noch nicht ausprobiert, aber der Reiz für mich ist hoch, da ich denke, dass
Testen von Anfang an sich nur positiv auf die Software auswirken kann.

Auch mein Interesse für die Automatisierung möchte ich in dieser Bachelorarbeit 
widerspiegeln, da dies den Entwicklungsprozess erheblich erleichtert und schneller
gestaltet.
	
	\onecolumn
	
	% einfacher Zeilenabstand
	\singlespacing
	
	% Literaturliste soll im Inhaltsverzeichnis auftauchen
	\newpage
	\addcontentsline{toc}{section}{Literaturverzeichnis}
	
	% Literaturverzeichnis anzeigen
	\renewcommand\refname{Literaturverzeichnis}
	\bibliography{expose}
\end{document}